\subsection{POWERS}

Consider \(j\) points of \(\mathbf{G}\) :

\[
\begin{array}{l} x (1) = (x (1) _ {1}, x (1) _ {2}, \ldots , x (1) _ {n}) \\ x (2) = (x (2) _ {1}, x (2) _ {2}, \ldots , x (2) _ {n}) \\ \cdot \cdot \cdot \\ x (j) = (x (j) _ {1}, x (j) _ {2}, \ldots , x (j) _ {n}). \end{array}
\]

The mapping \((x(1), x(2), \ldots, x(j)) \mapsto x(1).x(2) \ldots x(j)\) admits an expansion as an integral series about the origin:

\[
x (1). x (2) \dots x (j) = \sum_ {\alpha (1), \alpha (2), \dots , \alpha (j) \in \mathbf {N} ^ {n}} a _ {\alpha (1), \dots , \alpha (j)} x (1) ^ {\alpha (1)} \dots x (j) ^ {\alpha (j)}\tag{14}
\]

where the \(a_{\alpha(1), \ldots, \alpha(j)}\) are elements of \(\mathbf{K}^n\) . We write, for \(j = 0, 1, 2, \ldots\) ,

\[
\psi_ {j} (x) = \sum_ {\alpha (1) \neq 0, \dots , \alpha (j) \neq 0} a _ {\alpha (1), \dots , \alpha (j)} x ^ {\alpha (1) + \dots + \alpha (j)}\tag{15}
\]

where the right hand side is a convergent integral series in the variable \(x \in K^{n}\) . This series is obtained by suppressing in (14) the terms in which one of the variables \(x(1), \ldots, x(j)\) does not occur explicitly and then writing \(x(1) = x(2) = \cdots = x(j) = x\) .

If \(t \in K\) , all the t-th power mappings of G have the same expansion as an integral series about the origin, since any two of them coincide on a neighbourhood of 0. We denote this expansion as an integral series by \(x^{[t]}\) .

PROPOSITION 2. (i) \(\psi_{j} \equiv 0 \bmod \deg j\) .

\begin{enumerate}
\def\labelenumi{(\roman{enumi})}
\setcounter{enumi}{1}
\tightlist
\item
  If \(t\in \mathbf{K}\) , then
\end{enumerate}

\[
x ^ {[ t ]} = \sum_ {i = 1} ^ {\infty} \binom {t} {i} \psi_ {i} (x)\tag{16}
\]

where the formal power series on the right is meaningful because of (i).

(We write

\[
\binom {t} {i} = \frac {t (t - 1) \dots (t - i + 1)}{i !}
\]

for all \(t\in \mathbf{K}.\)

Assertion (i) is obvious from the definition of the \(\psi_{j}\) .

We prove (ii) for t an integer \(\geqslant0\) . By (14),

\[
x ^ {[ t ]} = \sum_ {\alpha (1),\ldots,\alpha (t)\in(\mathbf{N}^{n})^{t}} a _ {\alpha (1), \dots , \alpha (t)} x ^ {\alpha (1) + \dots + \alpha (t)}.\tag{17}
\]

For \(\alpha = (\alpha(1),\ldots,\alpha(t))\in(\mathbf{N}^{n})^{t}\) , let \(\sigma(\alpha)\) denote the set of \(j\in\{1,2,\ldots,t\}\) such that \(\alpha(j)\neq0\) . If, in the sum (17), we group together the terms for which \(\sigma(\alpha)\) is the same, then

\[
x ^ {[ t ]} = \sum_ {\sigma \subset (1, t)} h _ {t, \sigma} (x)\tag{18}
\]

where

\[
h _ {t, \sigma} (x) = \sum_ {\sigma (\alpha) = \sigma} a _ {\alpha (1), \dots , \alpha (t)} x ^ {\alpha (1) + \dots + \alpha (t)}.\tag{19}
\]

Let \(\sigma = \{j_1, j_2, \ldots, j_q\}\) with \(j_1 < j_2 < \ldots < j_q\) . In (14) (where \(j\) is replaced by \(t\) ), we substitute 0 for \(x(k)\) for \(k \notin \sigma\) ; as 0 is the identity element of G, we obtain the expansion of \(x(j_1).x(j_2)..x(j_q)\) as an integral series about the origin:

\[
x (j _ {1}). x (j _ {2}) \dots x (j _ {q}) = \sum_ {\sigma (\alpha) \subset \sigma} a _ {\alpha (1), \dots , \alpha (t)} x (j _ {1}) ^ {\alpha (j _ {1})} x (j _ {2}) ^ {\alpha (j _ {2})} \dots x (j _ {q}) ^ {\alpha (j _ {q})}
\]

and hence, by the definition of \(\psi_q\) :

\[
\psi_ {q} (x) = \sum_ {\sigma (\alpha) = \sigma} a _ {\alpha (1), \dots , \alpha (t)} x ^ {\alpha (j _ {1}) + \dots + \alpha (j _ {q})}.\tag{20}
\]

By (19) and (20), we see that \(h_{t,\sigma}(x) = \psi_{\mathrm{Card}\sigma}(x)\) . Then (18) implies

\[
x ^ {[ t ]} = \sum_ {i = 0} ^ {t} \binom {t} {i} \psi_ {i} (x) = \sum_ {i = 0} ^ {\infty} \binom {t} {i} \psi_ {i} (x).
\]

Then we write \(x^{[t]} = \sum_{i=0}^{\infty} \binom{t}{i}\psi_i(x)\) for all \(t\in\mathbf{K}\). In the integral series \(x^{[t]}\) and \(x^{[t']}\), each coefficient is a polynomial function of \(t\). For this is obvious for \(x^{[t']}\). As far as \(x^{[t]}\) is concerned, it suffices to prove that, for all \(u\in\mathrm{U}(G)\), the image of \(u\) under \(x\mapsto x^{[t]}\) is a polynomial function of \(t\). Now, for \(u\in\mathrm{U}^{m}(G)\), this image is \(t^{m}u\) (§ 4, no. 3, Proposition 7 (iv)).

As \(x^{[t]} = x^{[t']}\) for \(t\) an integer \(\geqslant 0\) , it then follows that \(x^{[t]} = x^{[t']}\) for all \(t \in \mathbf{K}\) .

Remarks. (1) We write condition (ii) of Proposition 2 for \(t\) an integer \(\geqslant 0\) :

\[
\begin{array}{r l} & 0 = \psi_ {0} (x) \\ & x = \psi_ {0} (x) + \psi_ {1} (x) \\ & x ^ {[ 2 ]} = \psi_ {0} (x) + 2 \psi_ {1} (x) + \psi_ {2} (x) \\ & \cdot \cdot \cdot . \end{array}
\]

These formulae suffice to determine the \(\psi_{i}\) .

\begin{enumerate}
\def\labelenumi{(\arabic{enumi})}
\setcounter{enumi}{1}
\tightlist
\item
  We see that \(\psi_0(x) = 0\) , \(\psi_1(x) = x\) , \(\psi_2(x) = x^{[2]} - 2x\) ,
\end{enumerate}

\[
x ^ {[ - 1 ]} = \sum_ {i = 1} ^ {\infty} (- 1) ^ {i} \psi_ {i} (x).
\]

\begin{enumerate}
\def\labelenumi{(\arabic{enumi})}
\setcounter{enumi}{2}
\tightlist
\item
  The above expression for \(\psi_{2}\) and formula (6) prove that
\end{enumerate}

\[
\psi_ {2} (x) \equiv \mathrm{B} (x, x) \mod \deg 3.\tag{21}
\]

Using Proposition 2, (i) and (ii), we see that

\[
x ^ {[ t ]} \equiv t x + \binom {t} {2} \mathrm{B} (x, x) \mod \deg 3.\tag{22}
\]

\begin{enumerate}
\def\labelenumi{(\arabic{enumi})}
\setcounter{enumi}{3}
\tightlist
\item
  Let \(\psi_{p,m}(x)\) and \(h_{t,m}(x)\) denote the homogeneous components of \(\psi_p(x)\) and \(x^{[t]}\) of degree \(m\) . Then \(\psi_{p,m} = 0\) for \(m < p\) . On the other hand, Proposition 2 (ii) gives
\end{enumerate}

\[
h _ {t, m} (x) = \sum_ {p \leqslant m} \frac {t (t - 1) \dots (t - p + 1)}{p !} \psi_ {p, m} (x)\tag{23}
\]

that is

\[
h _ {t, m} (x) = \sum_ {1 \leqslant r \leqslant m} t ^ {r} \phi_ {r, m} (x)\tag{24}
\]

where the \(\phi_{r,m}\) are homogeneous polynomial mappings of degree m of \(K^{n}\) into \(K^{n}\) . In particular, by (23),

\begin{enumerate}
\def\labelenumi{(\arabic{enumi})}
\setcounter{enumi}{24}
\tightlist
\item
\end{enumerate}

\[
\phi_ {1, m} (x) = \sum_ {p \leqslant m} \frac {(- 1) ^ {p - 1}}{p} \psi_ {p, m} (x)\tag{25}
\]

\[
\phi_ {m, m} (x) = \frac {1}{m !} \psi_ {m, m} (x).\tag{26}
\]

\begin{enumerate}
\def\labelenumi{(\arabic{enumi})}
\setcounter{enumi}{4}
\tightlist
\item
  If \(\mathbf{K}\) is of characteristic \(>0\) , the results of nos. 1 and 2 remain true, provided \(e_{\alpha}\) , in no. 2, is defined by \(\langle e_{\alpha}, \sum_{\beta} \lambda_{\beta} x^{\beta} \rangle = \lambda_{\alpha}\) . In no. 3, if the \(\psi_j\) are defined as above, the argument again proves that \(x^{[t]} = \sum_{i=1}^{\infty} \binom{t}{i} \psi_i(x)\) if \(t \in \mathbf{N}\) .
\end{enumerate}

